\section{Motivación}
El estado y mantenimiento de los neumáticos de los vehículos representan variables críticas para la seguridad vial, la integridad de los pasajeros y la eficiencia del transporte industrial. De acuerdo con informes globales de seguridad en el transporte, una fracción considerable de los accidentes de tránsito causados por fallos mecánicos es imputable al desgaste excesivo, roturas estructurales, grietas por resequedad y deformaciones en la banda de rodadura de las llantas. Por consiguiente, la inspección oportuna y el diagnóstico preciso del estado de la textura del neumático constituyen problemas de primer orden en la seguridad pública y el control de calidad en plantas de fabricación.

Tradicionalmente, la inspección de defectos en neumáticos ha sido un proceso manual altamente dependiente de operarios calificados. Este método tradicional presenta diversas desventajas críticas:
\begin{itemize}
    \item \textbf{Subjetividad y Fatiga Humana:} La evaluación visual depende de la apreciación del operario, la cual decae en precisión a lo largo de turnos de trabajo prolongados.
    \item \textbf{Baja Escalabilidad:} Inspeccionar manualmente cada neumático en una línea de producción de alto volumen o en centros de revisión masivos resulta ineficiente en tiempo y costo.
    \item \textbf{Falta de Estandarización:} La carencia de criterios cuantitativos repetibles dificulta el registro histórico y la trazabilidad de la calidad.
\end{itemize}

Con la maduración del aprendizaje profundo y la visión por computadora, surge una oportunidad tecnológica para automatizar esta tarea mediante Redes Neuronales Convolucionales (CNN). Resolver este problema de clasificación de texturas de neumáticos es de suma importancia por tres razones fundamentales:
\begin{enumerate}
    \item \textbf{Prevención Activa de Accidentes:} Un sistema automatizado y continuo de detección de fallos permite retirar neumáticos defectuosos antes de que ocurran reventones en autopistas a alta velocidad.
    \item \textbf{Optimización en Manufactura:} La integración en líneas de ensamble automatiza el control de calidad, asegurando que ningún neumático defectuoso salga de la planta, reduciendo devoluciones y costes de garantía.
    \item \textbf{Diagnóstico no Invasivo:} Permite realizar análisis mediante imágenes externas de manera rápida y sin necesidad de desmontar el neumático o realizar pruebas mecánicas destructivas.
\end{enumerate}

Sin embargo, para lograr un despliegue viable, los modelos computacionales deben superar retos severos como el desbalance inherente de datos (los neumáticos defectuosos son mucho más raros que los sanos) y la necesidad de interpretabilidad en sus decisiones, aspectos que motivan las estrategias experimentales adoptadas en este desarrollo.
